\documentclass[11pt, ukrainian]{article}
\setlength\textwidth{180mm} \setlength\textheight{277mm}
\setlength\voffset{-38mm} \setlength\hoffset{-28mm}
%\setlength\parindent{0mm}
\pagestyle{empty}
\usepackage[T2A]{fontenc}
\usepackage[utf8]{inputenc}
\usepackage[ukrainian]{babel}
\begin{document}
{\Large \center  Варіанти завдань до лабораторної роботи №4\par 2020/21 н.р.\par \bigskip\bigskip}

Порядок полів у записах та інформацію додаткового файлу фіксовано у варіанті.

Для порівняння числових значень використовується традиційний порядок.

Для порівняння рядкових значень використовується лексикографічний порядок. 

У випадку сортування за кількома полями результуюче сортування будується 
за порядками окремих полів за допомогою конструкції лексикографічного порядку. 
Послідовність полів, за якими має відбуватися сортування, задається в умові. 
Якщо не вказано інше, то поле сортується за зростанням.

Якщо умова вимагає знайти $k$ най… і $k$-е місце займають кілька, то слід виводити всіх їх.

\newpage
\par \bigskip

{\center Варіант  \No \textbf { 1 }\par}
\bigskip
%type = tasks
%variant = 2
Обробляються результати розв’язування студентами задач на контрольних.
Одиничний факт роз\-в’я\-зу\-вання студентом задачі надалі розуміється як спроба.

\underline{Записи основного файлу містять поля}: рік, по-батькові, ім’я, прізвище, умова задачі, відсоток набраних за розв’язання задачі балів.

\underline{У допоміжному файлі наявні ключі}: кількість спроб з відсотком не менше 50\%, кількість записів.

\underline{Знайти} всі задачі, які були найкраще розв’язані в середньому (середнє значення округлити до цілого). Вивести по кожному з них інформацію:\par
{\noindent}--- на першому рядку:\par
умова, середня розв’язуваність, відсоток спроб краще середнього (з округленням до цілого);

{\noindent}--- на наступних рядках, починаючи з табуляції, вивести спроби, що набрали якнайменш 75\% за цю задачу (по одному на рядок):\par
відсоток набраних за задачу балів, прізвище, ім’я, по-батькові, рік

{\noindent}у такому сортуванні: відсоток набраних за задачу балів (за спаданням), прізвище, ім’я, по-батькові, рік.

%Задачі сортувати: середня розв’язуваність (за спаданням), умова.

\bigskip\textit{Обмеження}

Студент міг розв’язувати задачу кілька разів.

Відсоток набраних за розв’язання задачі балів є цілим від 0 до 100.
День, місяць, рік є числовими полями, що зображують дату.
Раніше 2001 року статистика не велася. Рік є 4-значним.
Решта полів є непорожніми рядками.
Рядки починатися та закінчуватися білими символами не можуть.

Прізвище, ім’я та по-батькові (за наявності у варіанті) можуть мати до 20, 30, 28 відповідно символів включно кожне. 
Крім літер алфавіту можуть містити апостроф, дефіс та пробіл.

Умова задачі може мати від 3 до 52 символів включно.
Крім літер алфавіту може містити подвійні лапки, десяткові цифри, пробіл, дефіс,
а також символи . , \_ \{ \} \$ $+$ $*$ {\textbackslash} ( ) \% $[$ $]$ .

Студенти однозначно ідентифікуються прізвищем, ім’ям та по-батькові (за наявності у варіанті).
Задачі однозначно ідентифікуються умовою.

\bigskip\textit{Для деяких варіантів може знадобитись}

Розв'язуваність задачі --- середній відсоток набраних за задачу балів.

Розв'язуваність студента --- середній відсоток набраних студентом балів по всіх його спробах.

Розв'язуваність --- середній відсоток набраних студентами за задачі балів (рахується по всім даним).


\par
\newpage
{\center Варіант  \No \textbf { 2 }\par}
\bigskip
%type = tasks
%variant = 3
Обробляються результати розв’язування студентами задач на контрольних.
Одиничний факт роз\-в’я\-зу\-вання студентом задачі надалі розуміється як спроба.

\underline{Записи основного файлу містять поля}: відсоток набраних за розв’язання задачі балів, день, прізвище, місяць, рік, ім’я, умова задачі.

\underline{У допоміжному файлі наявні ключі}: загальна кількість записів, сумарний відсоток по всім записам файлу.

\underline{Знайти} всі задачі, які не розв’язав ніхто (з відсотком 0). Вивести по кожному з них інформацію:\par
{\noindent}--- на першому рядку:\par
умова, кількість спроб розв’язати;

{\noindent}--- на наступних рядках, починаючи з табуляції, вивести спроби по цих задачах (по одному на рядок):\par
прізвище, ім’я, рік, місяць, день

{\noindent}у такому сортуванні: рік (за спаданням), місяць, день, прізвище, ім’я.

%Задачі сортувати: кількість спроб (за спаданням), умова.

\bigskip\textit{Обмеження}

Студент міг розв’язувати задачу кілька разів.

Відсоток набраних за розв’язання задачі балів є цілим від 0 до 100.
День, місяць, рік є числовими полями, що зображують дату.
Раніше 2001 року статистика не велася. Рік є 4-значним.
Решта полів є непорожніми рядками.
Рядки починатися та закінчуватися білими символами не можуть.

Прізвище, ім’я та по-батькові (за наявності у варіанті) можуть мати до 23, 24, 23 відповідно символів включно кожне. 
Крім літер алфавіту можуть містити апостроф, дефіс та пробіл.

Умова задачі може мати від 4 до 66 символів включно.
Крім літер алфавіту може містити подвійні лапки, десяткові цифри, пробіл, дефіс,
а також символи . , \_ \{ \} \$ $+$ $*$ {\textbackslash} ( ) \% $[$ $]$ .

Студенти однозначно ідентифікуються прізвищем, ім’ям та по-батькові (за наявності у варіанті).
Задачі однозначно ідентифікуються умовою.

\bigskip\textit{Для деяких варіантів може знадобитись}

Розв'язуваність задачі --- середній відсоток набраних за задачу балів.

Розв'язуваність студента --- середній відсоток набраних студентом балів по всіх його спробах.

Розв'язуваність --- середній відсоток набраних студентами за задачі балів (рахується по всім даним).


\par
\newpage
{\center Варіант  \No \textbf { 3 }\par}
\bigskip
%type = tasks
%variant = 16
Обробляються результати розв’язування студентами задач на контрольних.
Одиничний факт роз\-в’я\-зу\-вання студентом задачі надалі розуміється як спроба.

\underline{Записи основного файлу містять поля}: відсоток набраних за розв’язання задачі балів, ім’я, рік, прізвище, умова задачі, по-батькові.

\underline{У допоміжному файлі наявні ключі}: загальна кількість студентів, найбільше значення відсотка, найменше значення відсотка.

\underline{Знайти} всіх студентів, у яких всіх спроби лежать у межах від 75\% до 90\% включно. Вивести по кожному з них інформацію:\par
{\noindent}--- на першому рядку:\par
середня розв’язуваність студента (округлена до одного знаку), ім’я, по-батькові, прізвище, кількість спроб;

{\noindent}--- на наступних рядках, починаючи з табуляції, вивести їх спроби (по одному на рядок):\par
умова, рік, відсоток набраних за задачу балів

{\noindent}у такому сортуванні: рік (за спаданням), відсоток, умова.

%Студентів сортувати: середня розв’язуваність, кількість спроб, прізвище, ім’я, по-батькові.

\bigskip\textit{Обмеження}

Студент міг розв’язувати задачу кілька разів.

Відсоток набраних за розв’язання задачі балів є цілим від 0 до 100.
День, місяць, рік є числовими полями, що зображують дату.
Раніше 2001 року статистика не велася. Рік є 4-значним.
Решта полів є непорожніми рядками.
Рядки починатися та закінчуватися білими символами не можуть.

Прізвище, ім’я та по-батькові (за наявності у варіанті) можуть мати до 25, 20, 24 відповідно символів включно кожне. 
Крім літер алфавіту можуть містити апостроф, дефіс та пробіл.

Умова задачі може мати від 5 до 67 символів включно.
Крім літер алфавіту може містити подвійні лапки, десяткові цифри, пробіл, дефіс,
а також символи . , \_ \{ \} \$ $+$ $*$ {\textbackslash} ( ) \% $[$ $]$ .

Студенти однозначно ідентифікуються прізвищем, ім’ям та по-батькові (за наявності у варіанті).
Задачі однозначно ідентифікуються умовою.

\bigskip\textit{Для деяких варіантів може знадобитись}

Розв'язуваність задачі --- середній відсоток набраних за задачу балів.

Розв'язуваність студента --- середній відсоток набраних студентом балів по всіх його спробах.

Розв'язуваність --- середній відсоток набраних студентами за задачі балів (рахується по всім даним).


\par
\newpage
{\center Варіант  \No \textbf { 4 }\par}
\bigskip
%type = tasks
%variant = 10
Обробляються результати розв’язування студентами задач на контрольних.
Одиничний факт роз\-в’я\-зу\-вання студентом задачі надалі розуміється як спроба.

\underline{Записи основного файлу містять поля}: по-батькові, рік, умова задачі, відсоток набраних за розв’я\-зан\-ня задачі балів, ім’я, прізвище.

\underline{У допоміжному файлі наявні ключі}: найменший відсоток набраних балів; середній відсоток набраних балів, округлений до цілого.

\underline{Знайти} всі задачі, які хоча б раз були розв’язані не менш ніж на 90\%. Вивести по кожному з них інформацію:\par
{\noindent}--- на першому рядку:\par
середня розв’язуваність (округлена до 1 знаку), умова, кількість спроб;

{\noindent}--- на наступних рядках, починаючи з табуляції, вивести спроби по цих задачах з відсотком не більше 75\% (по одному на рядок):\par
прізвище, ім’я, по-батькові, відсоток набраних балів, рік

{\noindent}у такому сортуванні: відсоток набраних балів (за спаданням), рік, прізвище, ім’я, по-батькові.

%Задачі сортувати: середня розв’язуваність, умова.

\bigskip\textit{Обмеження}

Студент міг розв’язувати задачу кілька разів.

Відсоток набраних за розв’язання задачі балів є цілим від 0 до 100.
День, місяць, рік є числовими полями, що зображують дату.
Раніше 2001 року статистика не велася. Рік є 4-значним.
Решта полів є непорожніми рядками.
Рядки починатися та закінчуватися білими символами не можуть.

Прізвище, ім’я та по-батькові (за наявності у варіанті) можуть мати до 29, 22, 26 відповідно символів включно кожне. 
Крім літер алфавіту можуть містити апостроф, дефіс та пробіл.

Умова задачі може мати від 5 до 56 символів включно.
Крім літер алфавіту може містити подвійні лапки, десяткові цифри, пробіл, дефіс,
а також символи . , \_ \{ \} \$ $+$ $*$ {\textbackslash} ( ) \% $[$ $]$ .

Студенти однозначно ідентифікуються прізвищем, ім’ям та по-батькові (за наявності у варіанті).
Задачі однозначно ідентифікуються умовою.

\bigskip\textit{Для деяких варіантів може знадобитись}

Розв'язуваність задачі --- середній відсоток набраних за задачу балів.

Розв'язуваність студента --- середній відсоток набраних студентом балів по всіх його спробах.

Розв'язуваність --- середній відсоток набраних студентами за задачі балів (рахується по всім даним).


\par
\newpage
{\center Варіант  \No \textbf { 5 }\par}
\bigskip
%type = tasks
%variant = 13
Обробляються результати розв’язування студентами задач на контрольних.
Одиничний факт роз\-в’я\-зу\-вання студентом задачі надалі розуміється як спроба.

\underline{Записи основного файлу містять поля}: умова задачі, відсоток набраних за розв’язання задачі балів, місяць, рік, день, ім’я, прізвище, по-батькові.

\underline{У допоміжному файлі наявні ключі}: сума днів, кількість спроб на 100\%.

\underline{Знайти} всі задачі, для яких існує спроба з відсотком розв’язання рівно 77\%. Вивести по кожному з них інформацію:\par
{\noindent}--- на першому рядку:\par
умова, кількість спроб кращих за 75\%, кількість спроб, рік, місяць, день останньої найкращої спроби;

{\noindent}--- на наступних рядках, починаючи з табуляції, вивести спроби по цих задачах (по одному на рядок):\par
прізвище, ім’я, по-батькові, відсоток набраних балів, рік, місяць, день

{\noindent}у такому сортуванні: рік, місяць, день, відсоток набраних балів, прізвище, ім’я, по-батькові.

%Задачі сортувати: кількість спроб кращих за 75\%, умова.

\bigskip\textit{Обмеження}

Студент міг розв’язувати задачу кілька разів.

Відсоток набраних за розв’язання задачі балів є цілим від 0 до 100.
День, місяць, рік є числовими полями, що зображують дату.
Раніше 2001 року статистика не велася. Рік є 4-значним.
Решта полів є непорожніми рядками.
Рядки починатися та закінчуватися білими символами не можуть.

Прізвище, ім’я та по-батькові (за наявності у варіанті) можуть мати до 22, 21, 21 відповідно символів включно кожне. 
Крім літер алфавіту можуть містити апостроф, дефіс та пробіл.

Умова задачі може мати від 3 до 70 символів включно.
Крім літер алфавіту може містити подвійні лапки, десяткові цифри, пробіл, дефіс,
а також символи . , \_ \{ \} \$ $+$ $*$ {\textbackslash} ( ) \% $[$ $]$ .

Студенти однозначно ідентифікуються прізвищем, ім’ям та по-батькові (за наявності у варіанті).
Задачі однозначно ідентифікуються умовою.

\bigskip\textit{Для деяких варіантів може знадобитись}

Розв'язуваність задачі --- середній відсоток набраних за задачу балів.

Розв'язуваність студента --- середній відсоток набраних студентом балів по всіх його спробах.

Розв'язуваність --- середній відсоток набраних студентами за задачі балів (рахується по всім даним).


\par
\newpage
{\center Варіант  \No \textbf { 6 }\par}
\bigskip
%type = tasks
%variant = 15
Обробляються результати розв’язування студентами задач на контрольних.
Одиничний факт роз\-в’я\-зу\-вання студентом задачі надалі розуміється як спроба.

\underline{Записи основного файлу містять поля}: ім’я, по-батькові, рік, відсоток набраних за розв’язання задачі балів, прізвище, умова задачі.

\underline{У допоміжному файлі наявні ключі}: найбільша довжина умови, загальна кількість записів.

\underline{Знайти} всіх студентів, що не мають спроб з відсотком не менше 90\%. Вивести по кожному з них інформацію:\par
{\noindent}--- на першому рядку:\par
ім’я, по-батькові, прізвище, середня розв’язуваність студента (округлена до цілого), кількість спроб;

{\noindent}--- на наступних рядках, починаючи з табуляції, вивести їх спроби (по одному на рядок):\par
умова, рік, відсоток набраних за задачу балів

{\noindent}у такому сортуванні: умова, відсоток, рік.

%Студентів сортувати: середня розв’язуваність, кількість спроб (за спаданням), прізвище, ім’я, по-батькові.

\bigskip\textit{Обмеження}

Студент міг розв’язувати задачу кілька разів.

Відсоток набраних за розв’язання задачі балів є цілим від 0 до 100.
День, місяць, рік є числовими полями, що зображують дату.
Раніше 2001 року статистика не велася. Рік є 4-значним.
Решта полів є непорожніми рядками.
Рядки починатися та закінчуватися білими символами не можуть.

Прізвище, ім’я та по-батькові (за наявності у варіанті) можуть мати до 26, 27, 20 відповідно символів включно кожне. 
Крім літер алфавіту можуть містити апостроф, дефіс та пробіл.

Умова задачі може мати від 4 до 50 символів включно.
Крім літер алфавіту може містити подвійні лапки, десяткові цифри, пробіл, дефіс,
а також символи . , \_ \{ \} \$ $+$ $*$ {\textbackslash} ( ) \% $[$ $]$ .

Студенти однозначно ідентифікуються прізвищем, ім’ям та по-батькові (за наявності у варіанті).
Задачі однозначно ідентифікуються умовою.

\bigskip\textit{Для деяких варіантів може знадобитись}

Розв'язуваність задачі --- середній відсоток набраних за задачу балів.

Розв'язуваність студента --- середній відсоток набраних студентом балів по всіх його спробах.

Розв'язуваність --- середній відсоток набраних студентами за задачі балів (рахується по всім даним).


\par
\newpage
{\center Варіант  \No \textbf { 7 }\par}
\bigskip
%type =     tasks
%variant = 1
Обробляються результати розв’язування студентами задач на контрольних.
Одиничний факт роз\-в’я\-зу\-вання студентом задачі надалі розуміється як спроба.

\underline{Записи основного файлу містять поля}: по-батькові, рік, відсоток набраних за розв’язання задачі балів, місяць, ім’я, прізвище, день, умова задачі.

\underline{У допоміжному файлі наявні ключі}: кількість різних задач, кількість записів.

\underline{Знайти} всі задачі, які були розв’язані на 100\% у найбільшій кількості спроб. Вивести по кожному з них інформацію:\par
{\noindent}--- на першому рядку:\par
умова, кількість спроб на 100\%, кількість спроб;

{\noindent}--- на наступних рядках, починаючи з табуляції, вивести спроби, що набрали якнайменш 75\% за цю задачу (по одному на рядок):\par
рік, місяць, відсоток набраних за задачу балів, прізвище, ім’я, по-батькові

{\noindent}у такому сортуванні: відсоток набраних за задачу балів (за спаданням), рік, місяць, прізвище, ім’я, по-батькові.

%Задачі сортувати: умова.

\bigskip\textit{Обмеження}

Студент міг розв’язувати задачу кілька разів.

Відсоток набраних за розв’язання задачі балів є цілим від 0 до 100.
День, місяць, рік є числовими полями, що зображують дату.
Раніше 2001 року статистика не велася. Рік є 4-значним.
Решта полів є непорожніми рядками.
Рядки починатися та закінчуватися білими символами не можуть.

Прізвище, ім’я та по-батькові (за наявності у варіанті) можуть мати до 28, 29, 22 відповідно символів включно кожне. 
Крім літер алфавіту можуть містити апостроф, дефіс та пробіл.

Умова задачі може мати від 4 до 59 символів включно.
Крім літер алфавіту може містити подвійні лапки, десяткові цифри, пробіл, дефіс,
а також символи . , \_ \{ \} \$ $+$ $*$ {\textbackslash} ( ) \% $[$ $]$ .

Студенти однозначно ідентифікуються прізвищем, ім’ям та по-батькові (за наявності у варіанті).
Задачі однозначно ідентифікуються умовою.

\bigskip\textit{Для деяких варіантів може знадобитись}

Розв'язуваність задачі --- середній відсоток набраних за задачу балів.

Розв'язуваність студента --- середній відсоток набраних студентом балів по всіх його спробах.

Розв'язуваність --- середній відсоток набраних студентами за задачі балів (рахується по всім даним).


\par
\newpage
{\center Варіант  \No \textbf { 8 }\par}
\bigskip
%type = tasks
%variant = 20
Обробляються результати розв’язування студентами задач на контрольних.
Одиничний факт роз\-в’я\-зу\-вання студентом задачі надалі розуміється як спроба.

\underline{Записи основного файлу містять поля}: по-батькові, умова задачі, відсоток набраних за розв’язання задачі балів, день, ім’я, прізвище, місяць, рік.

\underline{У допоміжному файлі наявні ключі}: середній відсоток з округленням до цілого, довжина найкоротшої умови.

\underline{Знайти} всіх студентів, що не мають спроб з відсотком менше 60\%. Вивести по кожному з них інформацію:\par
{\noindent}--- на першому рядку:\par
найкраща спроба, найгірша спроба, прізвище, ім’я, по-батькові, кількість спроб з відсотком 75\% та більше;

{\noindent}--- на наступних рядках, починаючи з табуляції, вивести їх спроби (по одному на рядок):\par
відсоток набраних за задачу балів, умова, рік

{\noindent}у такому сортуванні: відсоток (за спаданням), рік (за спаданням), місяць, умова.

%Студентів сортувати: найгірша спроба, кількість спроб з відсотком не менше 75\%, прізвище, ім’я, по-батькові.

\bigskip\textit{Обмеження}

Студент міг розв’язувати задачу кілька разів.

Відсоток набраних за розв’язання задачі балів є цілим від 0 до 100.
День, місяць, рік є числовими полями, що зображують дату.
Раніше 2001 року статистика не велася. Рік є 4-значним.
Решта полів є непорожніми рядками.
Рядки починатися та закінчуватися білими символами не можуть.

Прізвище, ім’я та по-батькові (за наявності у варіанті) можуть мати до 30, 23, 25 відповідно символів включно кожне. 
Крім літер алфавіту можуть містити апостроф, дефіс та пробіл.

Умова задачі може мати від 5 до 69 символів включно.
Крім літер алфавіту може містити подвійні лапки, десяткові цифри, пробіл, дефіс,
а також символи . , \_ \{ \} \$ $+$ $*$ {\textbackslash} ( ) \% $[$ $]$ .

Студенти однозначно ідентифікуються прізвищем, ім’ям та по-батькові (за наявності у варіанті).
Задачі однозначно ідентифікуються умовою.

\bigskip\textit{Для деяких варіантів може знадобитись}

Розв'язуваність задачі --- середній відсоток набраних за задачу балів.

Розв'язуваність студента --- середній відсоток набраних студентом балів по всіх його спробах.

Розв'язуваність --- середній відсоток набраних студентами за задачі балів (рахується по всім даним).


\par
\newpage
{\center Варіант  \No \textbf { 9 }\par}
\bigskip
%type = tasks
%variant = 19
Обробляються результати розв’язування студентами задач на контрольних.
Одиничний факт роз\-в’я\-зу\-вання студентом задачі надалі розуміється як спроба.

\underline{Записи основного файлу містять поля}: відсоток набраних за розв’язання задачі балів, умова задачі, ім’я, по-батькові, рік, прізвище.

\underline{У допоміжному файлі наявні ключі}: найбільший відсоток, найменший відсоток, загальна кількість записів.

\underline{Знайти} всіх студентів, у яких більше половини спроб отримали не менше 75\%. Вивести по кожному з них інформацію:\par
{\noindent}--- на першому рядку:\par
прізвище, ім’я, по-батькові, найкраща спроба, середня розв’язуваність (округлена до цілого), кількість спроб з відсотком не нижче 75\%;

{\noindent}--- на наступних рядках, починаючи з табуляції, вивести їх спроби (по одному на рядок):\par
відсоток набраних за задачу балів, умова, рік

{\noindent}у такому сортуванні: відсоток (за спаданням), рік (за спаданням), умова.

%Студентів сортувати: найкраща спроба (за зростанням), середня розв’язуваність, прізвище, ім’я, по-батькові.

\bigskip\textit{Обмеження}

Студент міг розв’язувати задачу кілька разів.

Відсоток набраних за розв’язання задачі балів є цілим від 0 до 100.
День, місяць, рік є числовими полями, що зображують дату.
Раніше 2001 року статистика не велася. Рік є 4-значним.
Решта полів є непорожніми рядками.
Рядки починатися та закінчуватися білими символами не можуть.

Прізвище, ім’я та по-батькові (за наявності у варіанті) можуть мати до 24, 28, 30 відповідно символів включно кожне. 
Крім літер алфавіту можуть містити апостроф, дефіс та пробіл.

Умова задачі може мати від 3 до 60 символів включно.
Крім літер алфавіту може містити подвійні лапки, десяткові цифри, пробіл, дефіс,
а також символи . , \_ \{ \} \$ $+$ $*$ {\textbackslash} ( ) \% $[$ $]$ .

Студенти однозначно ідентифікуються прізвищем, ім’ям та по-батькові (за наявності у варіанті).
Задачі однозначно ідентифікуються умовою.

\bigskip\textit{Для деяких варіантів може знадобитись}

Розв'язуваність задачі --- середній відсоток набраних за задачу балів.

Розв'язуваність студента --- середній відсоток набраних студентом балів по всіх його спробах.

Розв'язуваність --- середній відсоток набраних студентами за задачі балів (рахується по всім даним).


\par
\newpage
{\center Варіант  \No \textbf { 10 }\par}
\bigskip
%type = tasks
%variant = 18
Обробляються результати розв’язування студентами задач на контрольних.
Одиничний факт роз\-в’я\-зу\-вання студентом задачі надалі розуміється як спроба.

\underline{Записи основного файлу містять поля}: по-батькові, прізвище, день, відсоток набраних за розв’язан\-ня задачі балів, місяць, умова задачі, ім’я, рік.

\underline{У допоміжному файлі наявні ключі}: найкоротша умова, кількість записів.

\underline{Знайти} всіх студентів, що зробили рівно одну спробу з відсотком 99\%. Вивести по кожному з них інформацію:\par
{\noindent}--- на першому рядку:\par
прізвище, ім’я, по-батькові, найкраща спроба, найгірша спроба, кількість спроб з відсотком 99\%;

{\noindent}--- на наступних рядках, починаючи з табуляції, вивести їх спроби (по одному на рядок):\par
рік, місяць, день, відсоток набраних за задачу балів, умова

{\noindent}у такому сортуванні: рік(за спаданням), місяць (за спаданням), день (за спаданням), відсоток (за спаданням), умова.

%Студентів сортувати: найкраща спроба (за спаданням), найгірша спроба, прізвище, ім’я, по-батькові.

\bigskip\textit{Обмеження}

Студент міг розв’язувати задачу кілька разів.

Відсоток набраних за розв’язання задачі балів є цілим від 0 до 100.
День, місяць, рік є числовими полями, що зображують дату.
Раніше 2001 року статистика не велася. Рік є 4-значним.
Решта полів є непорожніми рядками.
Рядки починатися та закінчуватися білими символами не можуть.

Прізвище, ім’я та по-батькові (за наявності у варіанті) можуть мати до 21, 25, 29 відповідно символів включно кожне. 
Крім літер алфавіту можуть містити апостроф, дефіс та пробіл.

Умова задачі може мати від 5 до 68 символів включно.
Крім літер алфавіту може містити подвійні лапки, десяткові цифри, пробіл, дефіс,
а також символи . , \_ \{ \} \$ $+$ $*$ {\textbackslash} ( ) \% $[$ $]$ .

Студенти однозначно ідентифікуються прізвищем, ім’ям та по-батькові (за наявності у варіанті).
Задачі однозначно ідентифікуються умовою.

\bigskip\textit{Для деяких варіантів може знадобитись}

Розв'язуваність задачі --- середній відсоток набраних за задачу балів.

Розв'язуваність студента --- середній відсоток набраних студентом балів по всіх його спробах.

Розв'язуваність --- середній відсоток набраних студентами за задачі балів (рахується по всім даним).


\par
\newpage
{\center Варіант  \No \textbf { 11 }\par}
\bigskip
%type = tasks
%variant = 4
Обробляються результати розв’язування студентами задач на контрольних.
Одиничний факт роз\-в’я\-зу\-вання студентом задачі надалі розуміється як спроба.

\underline{Записи основного файлу містять поля}: відсоток набраних за розв’язання задачі балів, умова задачі, прізвище, місяць, рік, день, ім’я.

\underline{У допоміжному файлі наявні ключі}: сума років по всіх спробах, кількість різних студентів.

\underline{Знайти} всіх студентів, хто в середньому (округлити до цілого) розв’язував задачі найкраще. Вивести по кожному з них інформацію:\par
{\noindent}--- на першому рядку:\par
прізвище, ім’я, середня розв’язуваність студента, кількість спроб;

{\noindent}--- на наступних рядках, починаючи з табуляції, вивести їх спроби (по одному на рядок):\par
умова, рік, відсоток набраних за задачу балів

{\noindent}у такому сортуванні: відсоток набраних балів (за спаданням), рік (за спаданням), умова.

%Студентів сортувати: ім’я, прізвище.

\bigskip\textit{Обмеження}

Студент міг розв’язувати задачу кілька разів.

Відсоток набраних за розв’язання задачі балів є цілим від 0 до 100.
День, місяць, рік є числовими полями, що зображують дату.
Раніше 2001 року статистика не велася. Рік є 4-значним.
Решта полів є непорожніми рядками.
Рядки починатися та закінчуватися білими символами не можуть.

Прізвище, ім’я та по-батькові (за наявності у варіанті) можуть мати до 27, 26, 27 відповідно символів включно кожне. 
Крім літер алфавіту можуть містити апостроф, дефіс та пробіл.

Умова задачі може мати від 4 до 55 символів включно.
Крім літер алфавіту може містити подвійні лапки, десяткові цифри, пробіл, дефіс,
а також символи . , \_ \{ \} \$ $+$ $*$ {\textbackslash} ( ) \% $[$ $]$ .

Студенти однозначно ідентифікуються прізвищем, ім’ям та по-батькові (за наявності у варіанті).
Задачі однозначно ідентифікуються умовою.

\bigskip\textit{Для деяких варіантів може знадобитись}

Розв'язуваність задачі --- середній відсоток набраних за задачу балів.

Розв'язуваність студента --- середній відсоток набраних студентом балів по всіх його спробах.

Розв'язуваність --- середній відсоток набраних студентами за задачі балів (рахується по всім даним).


\par
\newpage
{\center Варіант  \No \textbf { 12 }\par}
\bigskip
%type = tasks
%variant = 9
Обробляються результати розв’язування студентами задач на контрольних.
Одиничний факт роз\-в’я\-зу\-вання студентом задачі надалі розуміється як спроба.

\underline{Записи основного файлу містять поля}: день, рік, прізвище, умова задачі, ім’я, відсоток набраних за розв’язання задачі балів, місяць.

\underline{У допоміжному файлі наявні ключі}: кількість задач, кількість записів.

\underline{Знайти} всі задачі, які в кожній спробі були розв’язані не більш ніж на 90\%. Вивести по кожному з них інформацію:\par
{\noindent}--- на першому рядку:\par
умова, кількість спроб, середня розв’язуваність (округлена до цілого);

{\noindent}--- на наступних рядках, починаючи з табуляції, вивести спроби по цих задачах (по одному на рядок):\par
рік, прізвище, ім’я, відсоток набраних балів

{\noindent}у такому сортуванні: прізвище, ім’я, відсоток набраних балів (за спаданням), рік.

%Задачі сортувати: кількість спроб, умова.

\bigskip\textit{Обмеження}

Студент міг розв’язувати задачу кілька разів.

Відсоток набраних за розв’язання задачі балів є цілим від 0 до 100.
День, місяць, рік є числовими полями, що зображують дату.
Раніше 2001 року статистика не велася. Рік є 4-значним.
Решта полів є непорожніми рядками.
Рядки починатися та закінчуватися білими символами не можуть.

Прізвище, ім’я та по-батькові (за наявності у варіанті) можуть мати до 21, 23, 22 відповідно символів включно кожне. 
Крім літер алфавіту можуть містити апостроф, дефіс та пробіл.

Умова задачі може мати від 3 до 61 символів включно.
Крім літер алфавіту може містити подвійні лапки, десяткові цифри, пробіл, дефіс,
а також символи . , \_ \{ \} \$ $+$ $*$ {\textbackslash} ( ) \% $[$ $]$ .

Студенти однозначно ідентифікуються прізвищем, ім’ям та по-батькові (за наявності у варіанті).
Задачі однозначно ідентифікуються умовою.

\bigskip\textit{Для деяких варіантів може знадобитись}

Розв'язуваність задачі --- середній відсоток набраних за задачу балів.

Розв'язуваність студента --- середній відсоток набраних студентом балів по всіх його спробах.

Розв'язуваність --- середній відсоток набраних студентами за задачі балів (рахується по всім даним).


\par
\newpage
{\center Варіант  \No \textbf { 13 }\par}
\bigskip
%type = tasks
%variant = 8
Обробляються результати розв’язування студентами задач на контрольних.
Одиничний факт роз\-в’я\-зу\-вання студентом задачі надалі розуміється як спроба.

\underline{Записи основного файлу містять поля}: ім’я, відсоток набраних за розв’язання задачі балів, по-бать\-ко\-ві, умова задачі, рік, прізвище.

\underline{У допоміжному файлі наявні ключі}: кількість різних задач, кількість спроб з 50\% балів.

\underline{Знайти} всі задачі, для яких існує спроба з відсотком розв’язання менш ніж 60\%. Вивести по кожному з них інформацію:\par
{\noindent}--- на першому рядку:\par
умова, кількість спроб, кількість спроб гірше 60\%, середня розв’язуваність (округлена до 1 знаку);

{\noindent}--- на наступних рядках, починаючи з табуляції, вивести спроби по цих задачах з відсотком не більше 75\% (по одному на рядок):\par
ім’я, по-батькові, прізвище, відсоток набраних балів, рік

{\noindent}у такому сортуванні: рік (за спаданням), ім’я, по-батькові, прізвище, відсоток набраних балів (за спаданням).

%Задачі сортувати: кількість спроб гірше 60\%, умова.

\bigskip\textit{Обмеження}

Студент міг розв’язувати задачу кілька разів.

Відсоток набраних за розв’язання задачі балів є цілим від 0 до 100.
День, місяць, рік є числовими полями, що зображують дату.
Раніше 2001 року статистика не велася. Рік є 4-значним.
Решта полів є непорожніми рядками.
Рядки починатися та закінчуватися білими символами не можуть.

Прізвище, ім’я та по-батькові (за наявності у варіанті) можуть мати до 30, 25, 27 відповідно символів включно кожне. 
Крім літер алфавіту можуть містити апостроф, дефіс та пробіл.

Умова задачі може мати від 4 до 58 символів включно.
Крім літер алфавіту може містити подвійні лапки, десяткові цифри, пробіл, дефіс,
а також символи . , \_ \{ \} \$ $+$ $*$ {\textbackslash} ( ) \% $[$ $]$ .

Студенти однозначно ідентифікуються прізвищем, ім’ям та по-батькові (за наявності у варіанті).
Задачі однозначно ідентифікуються умовою.

\bigskip\textit{Для деяких варіантів може знадобитись}

Розв'язуваність задачі --- середній відсоток набраних за задачу балів.

Розв'язуваність студента --- середній відсоток набраних студентом балів по всіх його спробах.

Розв'язуваність --- середній відсоток набраних студентами за задачі балів (рахується по всім даним).


\par
\newpage
{\center Варіант  \No \textbf { 14 }\par}
\bigskip
%type = tasks
%variant = 17
Обробляються результати розв’язування студентами задач на контрольних.
Одиничний факт роз\-в’я\-зу\-вання студентом задачі надалі розуміється як спроба.

\underline{Записи основного файлу містять поля}: день, ім’я, відсоток набраних за розв’язання задачі балів, рік, умова задачі, прізвище, місяць.

\underline{У допоміжному файлі наявні ключі}: загальна кількість студентів; найменше значення року, що зустрічається у файлі.

\underline{Знайти} всіх студентів, що не робили спроб ані з відсотком 0\%, ані з відсотком 100\%. Вивести по кожному з них інформацію:\par
{\noindent}--- на першому рядку:\par
прізвище, ім’я, по-батькові, кількість спроб, кількість спроб з відсотком 60\%;

{\noindent}--- на наступних рядках, починаючи з табуляції, вивести їх спроби (по одному на рядок):\par
відсоток набраних за задачу балів, рік, місяць, день умова

{\noindent}у такому сортуванні: умова, відсоток (за спаданням), рік, місяць, день.

%Студентів сортувати: кількість спроб (за спаданням), прізвище, ім’я.

\bigskip\textit{Обмеження}

Студент міг розв’язувати задачу кілька разів.

Відсоток набраних за розв’язання задачі балів є цілим від 0 до 100.
День, місяць, рік є числовими полями, що зображують дату.
Раніше 2001 року статистика не велася. Рік є 4-значним.
Решта полів є непорожніми рядками.
Рядки починатися та закінчуватися білими символами не можуть.

Прізвище, ім’я та по-батькові (за наявності у варіанті) можуть мати до 23, 26, 21 відповідно символів включно кожне. 
Крім літер алфавіту можуть містити апостроф, дефіс та пробіл.

Умова задачі може мати від 3 до 57 символів включно.
Крім літер алфавіту може містити подвійні лапки, десяткові цифри, пробіл, дефіс,
а також символи . , \_ \{ \} \$ $+$ $*$ {\textbackslash} ( ) \% $[$ $]$ .

Студенти однозначно ідентифікуються прізвищем, ім’ям та по-батькові (за наявності у варіанті).
Задачі однозначно ідентифікуються умовою.

\bigskip\textit{Для деяких варіантів може знадобитись}

Розв'язуваність задачі --- середній відсоток набраних за задачу балів.

Розв'язуваність студента --- середній відсоток набраних студентом балів по всіх його спробах.

Розв'язуваність --- середній відсоток набраних студентами за задачі балів (рахується по всім даним).


\par
\newpage
{\center Варіант  \No \textbf { 15 }\par}
\bigskip
%type = tasks
%variant = 6
Обробляються результати розв’язування студентами задач на контрольних.
Одиничний факт роз\-в’я\-зу\-вання студентом задачі надалі розуміється як спроба.

\underline{Записи основного файлу містять поля}: день, прізвище, ім’я, місяць, відсоток набраних за розв’язан\-ня задачі балів, умова задачі, по-батькові, рік.

\underline{У допоміжному файлі наявні ключі}: сума всіх днів, сума всіх місяців.

\underline{Знайти} всіх студентів, що мають найбільшу кількість спроб з відсотком 90\% та більше. Вивести по кожному з них інформацію:\par
{\noindent}--- на першому рядку:\par
ім’я, по-батькові, прізвище, кількість спроб більших 90\%, кількість спроб;

{\noindent}--- на наступних рядках, починаючи з табуляції, вивести їх спроби з відсотком більше 50\% (по одному на рядок):\par
рік, відсоток набраних за задачу балів, умова

{\noindent}у такому сортуванні: умова, відсоток набраних балів, рік(за спаданням).

%Студентів сортувати: прізвище, ім’я, по-батькові.

\bigskip\textit{Обмеження}

Студент міг розв’язувати задачу кілька разів.

Відсоток набраних за розв’язання задачі балів є цілим від 0 до 100.
День, місяць, рік є числовими полями, що зображують дату.
Раніше 2001 року статистика не велася. Рік є 4-значним.
Решта полів є непорожніми рядками.
Рядки починатися та закінчуватися білими символами не можуть.

Прізвище, ім’я та по-батькові (за наявності у варіанті) можуть мати до 20, 21, 25 відповідно символів включно кожне. 
Крім літер алфавіту можуть містити апостроф, дефіс та пробіл.

Умова задачі може мати від 5 до 64 символів включно.
Крім літер алфавіту може містити подвійні лапки, десяткові цифри, пробіл, дефіс,
а також символи . , \_ \{ \} \$ $+$ $*$ {\textbackslash} ( ) \% $[$ $]$ .

Студенти однозначно ідентифікуються прізвищем, ім’ям та по-батькові (за наявності у варіанті).
Задачі однозначно ідентифікуються умовою.

\bigskip\textit{Для деяких варіантів може знадобитись}

Розв'язуваність задачі --- середній відсоток набраних за задачу балів.

Розв'язуваність студента --- середній відсоток набраних студентом балів по всіх його спробах.

Розв'язуваність --- середній відсоток набраних студентами за задачі балів (рахується по всім даним).


\par
\newpage
{\center Варіант  \No \textbf { 16 }\par}
\bigskip
%type = tasks
%variant = 14
Обробляються результати розв’язування студентами задач на контрольних.
Одиничний факт роз\-в’я\-зу\-вання студентом задачі надалі розуміється як спроба.

\underline{Записи основного файлу містять поля}: день, прізвище, місяць, умова задачі, відсоток набраних за розв’язання задачі балів, ім’я, рік.

\underline{У допоміжному файлі наявні ключі}: загальна довжина прізвищ, найбільший відсоток набраних балів.

\underline{Знайти} всі задачі, за якими більше половини спроб отримало більше 60\%. Вивести по кожному з них інформацію:\par
{\noindent}--- на першому рядку:\par
умова, кількість спроб, кількість спроб кращих за 60\%, відсоток найкращої спроби;

{\noindent}--- на наступних рядках, починаючи з табуляції, вивести спроби по цих задачах (по одному на рядок):\par
прізвище, ім’я, відсоток набраних балів, рік, місяць, день

{\noindent}у такому сортуванні: рік, місяць, день, відсоток набраних балів, прізвище, ім’я.

%Задачі сортувати: відсоток найкращої спроби, кількість спроб, умова.

\bigskip\textit{Обмеження}

Студент міг розв’язувати задачу кілька разів.

Відсоток набраних за розв’язання задачі балів є цілим від 0 до 100.
День, місяць, рік є числовими полями, що зображують дату.
Раніше 2001 року статистика не велася. Рік є 4-значним.
Решта полів є непорожніми рядками.
Рядки починатися та закінчуватися білими символами не можуть.

Прізвище, ім’я та по-батькові (за наявності у варіанті) можуть мати до 28, 28, 29 відповідно символів включно кожне. 
Крім літер алфавіту можуть містити апостроф, дефіс та пробіл.

Умова задачі може мати від 3 до 53 символів включно.
Крім літер алфавіту може містити подвійні лапки, десяткові цифри, пробіл, дефіс,
а також символи . , \_ \{ \} \$ $+$ $*$ {\textbackslash} ( ) \% $[$ $]$ .

Студенти однозначно ідентифікуються прізвищем, ім’ям та по-батькові (за наявності у варіанті).
Задачі однозначно ідентифікуються умовою.

\bigskip\textit{Для деяких варіантів може знадобитись}

Розв'язуваність задачі --- середній відсоток набраних за задачу балів.

Розв'язуваність студента --- середній відсоток набраних студентом балів по всіх його спробах.

Розв'язуваність --- середній відсоток набраних студентами за задачі балів (рахується по всім даним).


\par
\newpage
{\center Варіант  \No \textbf { 17 }\par}
\bigskip
%type = tasks
%variant = 12
Обробляються результати розв’язування студентами задач на контрольних.
Одиничний факт роз\-в’я\-зу\-вання студентом задачі надалі розуміється як спроба.

\underline{Записи основного файлу містять поля}: місяць, день, прізвище, умова задачі, ім’я, по-батькові, відсоток набраних за розв’язання задачі балів, рік.

\underline{У допоміжному файлі наявні ключі}: найменший відсоток, найбільший відсоток, кількість різних задач.

\underline{Знайти} всі задачі, за якими не було жодної спроби краще 50\%. Вивести по кожному з них інформацію:\par
{\noindent}--- на першому рядку:\par
умова, кількість спроб кращих за 30\%, кількість спроб, відсоток найкращої спроби;

{\noindent}--- на наступних рядках, починаючи з табуляції, вивести спроби по цих задачах (по одному на рядок):\par
прізвище, ім’я, по-батькові, відсоток набраних балів, рік

{\noindent}у такому сортуванні: відсоток набраних балів (спочатку нулеві, далі інші за зростанням), рік, прізвище, ім’я, по-батькові.

%Задачі сортувати: відсоток найкращої спроби, кількість спроб, умова.

\bigskip\textit{Обмеження}

Студент міг розв’язувати задачу кілька разів.

Відсоток набраних за розв’язання задачі балів є цілим від 0 до 100.
День, місяць, рік є числовими полями, що зображують дату.
Раніше 2001 року статистика не велася. Рік є 4-значним.
Решта полів є непорожніми рядками.
Рядки починатися та закінчуватися білими символами не можуть.

Прізвище, ім’я та по-батькові (за наявності у варіанті) можуть мати до 24, 29, 24 відповідно символів включно кожне. 
Крім літер алфавіту можуть містити апостроф, дефіс та пробіл.

Умова задачі може мати від 4 до 54 символів включно.
Крім літер алфавіту може містити подвійні лапки, десяткові цифри, пробіл, дефіс,
а також символи . , \_ \{ \} \$ $+$ $*$ {\textbackslash} ( ) \% $[$ $]$ .

Студенти однозначно ідентифікуються прізвищем, ім’ям та по-батькові (за наявності у варіанті).
Задачі однозначно ідентифікуються умовою.

\bigskip\textit{Для деяких варіантів може знадобитись}

Розв'язуваність задачі --- середній відсоток набраних за задачу балів.

Розв'язуваність студента --- середній відсоток набраних студентом балів по всіх його спробах.

Розв'язуваність --- середній відсоток набраних студентами за задачі балів (рахується по всім даним).


\par
\newpage
{\center Варіант  \No \textbf { 18 }\par}
\bigskip
%type = tasks
%variant = 11
Обробляються результати розв’язування студентами задач на контрольних.
Одиничний факт роз\-в’я\-зу\-вання студентом задачі надалі розуміється як спроба.

\underline{Записи основного файлу містять поля}: ім’я, місяць, день, відсоток набраних за розв’язання задачі балів, умова задачі, прізвище, рік.

\underline{У допоміжному файлі наявні ключі}: найменший рік спроби, кількість записів.

\underline{Знайти} всі задачі, за якими всі спроби лежать у межах від 60\% до 90\% включно. Вивести по кожному з них інформацію:\par
{\noindent}--- на першому рядку:\par
кількість спроб, середня розв’язуваність (округлена до 1 знаку), умова;

{\noindent}--- на наступних рядках, починаючи з табуляції, вивести спроби по цих задачах (по одному на рядок):\par
прізвище, ім’я, відсоток набраних балів

{\noindent}у такому сортуванні: прізвище, ім’я, відсоток набраних балів (за спаданням).

%Задачі сортувати: кількість спроб, середня розв’язуваність, умова.

\bigskip\textit{Обмеження}

Студент міг розв’язувати задачу кілька разів.

Відсоток набраних за розв’язання задачі балів є цілим від 0 до 100.
День, місяць, рік є числовими полями, що зображують дату.
Раніше 2001 року статистика не велася. Рік є 4-значним.
Решта полів є непорожніми рядками.
Рядки починатися та закінчуватися білими символами не можуть.

Прізвище, ім’я та по-батькові (за наявності у варіанті) можуть мати до 27, 27, 23 відповідно символів включно кожне. 
Крім літер алфавіту можуть містити апостроф, дефіс та пробіл.

Умова задачі може мати від 5 до 51 символів включно.
Крім літер алфавіту може містити подвійні лапки, десяткові цифри, пробіл, дефіс,
а також символи . , \_ \{ \} \$ $+$ $*$ {\textbackslash} ( ) \% $[$ $]$ .

Студенти однозначно ідентифікуються прізвищем, ім’ям та по-батькові (за наявності у варіанті).
Задачі однозначно ідентифікуються умовою.

\bigskip\textit{Для деяких варіантів може знадобитись}

Розв'язуваність задачі --- середній відсоток набраних за задачу балів.

Розв'язуваність студента --- середній відсоток набраних студентом балів по всіх його спробах.

Розв'язуваність --- середній відсоток набраних студентами за задачі балів (рахується по всім даним).


\par
\newpage
{\center Варіант  \No \textbf { 19 }\par}
\bigskip
%type = tasks
%variant = 7
Обробляються результати розв’язування студентами задач на контрольних.
Одиничний факт роз\-в’я\-зу\-вання студентом задачі надалі розуміється як спроба.

\underline{Записи основного файлу містять поля}: ім’я, прізвище, відсоток набраних за розв’язання задачі балів, по-батькові, умова задачі, рік, місяць, день.

\underline{У допоміжному файлі наявні ключі}: сума довжин всіх прізвищ, кількість записів.

\underline{Знайти} всі задачі, які в кожній спробі були розв’язані не менш ніж на 60\%. Вивести по кожному з них інформацію:\par
{\noindent}--- на першому рядку:\par
умова, кількість спроб, кількість спроб не гірше 60\%, відсоток спроб не гірше 60\% до кількості спроб по задачі (з округленням до цілого);

{\noindent}--- на наступних рядках, починаючи з табуляції, вивести спроби по цих задачах (по одному на рядок):\par
рік, прізвище, ім’я, по-батькові, відсоток набраних балів

{\noindent}у такому сортуванні: прізвище, ім’я, по-батькові, відсоток набраних балів (за спаданням), рік.

%Задачі сортувати: відсоток спроб не гірше 60\% (з виконаним округленням, за спаданням), умова.

\bigskip\textit{Обмеження}

Студент міг розв’язувати задачу кілька разів.

Відсоток набраних за розв’язання задачі балів є цілим від 0 до 100.
День, місяць, рік є числовими полями, що зображують дату.
Раніше 2001 року статистика не велася. Рік є 4-значним.
Решта полів є непорожніми рядками.
Рядки починатися та закінчуватися білими символами не можуть.

Прізвище, ім’я та по-батькові (за наявності у варіанті) можуть мати до 25, 30, 20 відповідно символів включно кожне. 
Крім літер алфавіту можуть містити апостроф, дефіс та пробіл.

Умова задачі може мати від 3 до 65 символів включно.
Крім літер алфавіту може містити подвійні лапки, десяткові цифри, пробіл, дефіс,
а також символи . , \_ \{ \} \$ $+$ $*$ {\textbackslash} ( ) \% $[$ $]$ .

Студенти однозначно ідентифікуються прізвищем, ім’ям та по-батькові (за наявності у варіанті).
Задачі однозначно ідентифікуються умовою.

\bigskip\textit{Для деяких варіантів може знадобитись}

Розв'язуваність задачі --- середній відсоток набраних за задачу балів.

Розв'язуваність студента --- середній відсоток набраних студентом балів по всіх його спробах.

Розв'язуваність --- середній відсоток набраних студентами за задачі балів (рахується по всім даним).


\par
\newpage
{\center Варіант  \No \textbf { 20 }\par}
\bigskip
%type = tasks
%variant = 5
Обробляються результати розв’язування студентами задач на контрольних.
Одиничний факт роз\-в’я\-зу\-вання студентом задачі надалі розуміється як спроба.

\underline{Записи основного файлу містять поля}: відсоток набраних за розв’язання задачі балів, рік, по-бать\-ко\-ві, умова задачі, ім’я, прізвище.

\underline{У допоміжному файлі наявні ключі}: сумарна довжина прізвищ по всіх спробах, кількість записів.

\underline{Знайти} всі спроби, коли задача була розв’язана краще, ніж в середньому по цій задачі. Вивести по кожному з них інформацію:\par
{\noindent}--- на першому рядку:\par
умова, середня розв’язуваність задачі (округлена до 1 знаку);

{\noindent}--- на наступних рядках, починаючи з табуляції, вивести спроби по цих задачах (по одному на рядок):\par
рік, прізвище, ім’я, відсоток набраних балів

{\noindent}у такому сортуванні: відсоток набраних балів, рік, прізвище, ім’я.

%Задачі сортувати: середня розв’язуваність, умова.

\bigskip\textit{Обмеження}

Студент міг розв’язувати задачу кілька разів.

Відсоток набраних за розв’язання задачі балів є цілим від 0 до 100.
День, місяць, рік є числовими полями, що зображують дату.
Раніше 2001 року статистика не велася. Рік є 4-значним.
Решта полів є непорожніми рядками.
Рядки починатися та закінчуватися білими символами не можуть.

Прізвище, ім’я та по-батькові (за наявності у варіанті) можуть мати до 22, 24, 30 відповідно символів включно кожне. 
Крім літер алфавіту можуть містити апостроф, дефіс та пробіл.

Умова задачі може мати від 4 до 63 символів включно.
Крім літер алфавіту може містити подвійні лапки, десяткові цифри, пробіл, дефіс,
а також символи . , \_ \{ \} \$ $+$ $*$ {\textbackslash} ( ) \% $[$ $]$ .

Студенти однозначно ідентифікуються прізвищем, ім’ям та по-батькові (за наявності у варіанті).
Задачі однозначно ідентифікуються умовою.

\bigskip\textit{Для деяких варіантів може знадобитись}

Розв'язуваність задачі --- середній відсоток набраних за задачу балів.

Розв'язуваність студента --- середній відсоток набраних студентом балів по всіх його спробах.

Розв'язуваність --- середній відсоток набраних студентами за задачі балів (рахується по всім даним).


\par
\newpage
\end{document}
